\subsection{Simulation}

The validation run uses the normalised plant \(K_m=1\), \(b=1\). The inertia
switches from the small flywheel to the large flywheel and back to the small
flywheel, abbreviated as S--L--S, with
\begin{equation}
  J=1
  \qquad\hbox{and}\qquad
  J=24 .
  \label{eq:simulation-inertias}
\end{equation}
The ratio \(24{:}1\) corresponds to the two flywheels. The run lasts
\(\SI{96}{s}\). The reference is a rectangular speed command with amplitude
\(\SI{50}{rad/s}\) and period \(\SI{16}{s}\), starting at \(t=\SI{4}{s}\) from
standstill. Measurement noise is included.

With \(\alpha=0.98\), the overshoot before the inertia change lies between
\(\SI{0.74}{\percent}\) and \(\SI{0.84}{\percent}\). The \(2\%\) settling time
is \(\SI{0.40}{s}\). After the L--S change the first transient reaches
\(\SI{1.83}{\percent}\) overshoot and settles in \(\SI{0.56}{s}\). The next
two transients give \(\SI{0.78}{\percent}\) with \(\SI{0.48}{s}\), and
\(\SI{0.77}{\percent}\) with \(\SI{0.40}{s}\), respectively.

The motor command remains inside the saturation limits \(\pm255\). During
standstill the trace of \(P\) remains constant at 200, showing that the
excitation gate prevents covariance windup. The maximum trace over the whole
run is also 200. After the S--L inertia change, \(b_0\) reconverges after
\(\SI{12.24}{s}\), or \(\SI{8.24}{s}\) after the first exciting reference edge.
After the L--S change, the command sign changes at \(0.12\) times per second;
this does not indicate a limit cycle.

The forgetting factor is deliberately set to \(\alpha=0.98\), rather than
\(\alpha=1.0\). With \(\alpha=1.0\), the covariance becomes too small during
the earlier part of the run. After the L--S inertia change the estimator then
does not relearn the new plant quickly enough, and the closed loop ends in a
limit cycle with command chattering.
